%; whizzy chapter -dvi
% -initex iniptex -latex platex -format platex -bibtex jbibtex -fmt fmt
% 以上 whizzytex を使用する場合の設定。
 
%     Tokyo Debian Meeting resources
%     Copyright (C) 2012 Junichi Uekawa
%     Copyright (C) 2011 Nobuhiro Iwamatsu

%     This program is free software; you can redistribute it and/or modify
%     it under the terms of the GNU General Public License as published by
%     the Free Software Foundation; either version 2 of the License, or
%     (at your option) any later version.

%     This program is distributed in the hope that it will be useful,
%     but WITHOUT ANY WARRANTY; without even the implied warranty of
%     MERCHANTABILITY or FITNESS FOR A PARTICULAR PURPOSE.  See the
%     GNU General Public License for more details.

%     You should have received a copy of the GNU General Public License
%     along with this program; if not, write to the Free Software
%     Foundation, Inc., 51 Franklin St, Fifth Floor, Boston, MA  02110-1301 USA

%  preview (shell-command (concat "evince " (replace-regexp-in-string "tex$" "pdf"(buffer-file-name)) "&"))

%%ここからヘッダ開始。

\documentclass[mingoth,a4paper]{jsarticle}
\usepackage{monthlyreport}
% 日付を定義する、毎月変わります。
\newcommand{\debmtgyear}{2014}
\newcommand{\debmtgmonth}{07}
\newcommand{\debmtgdate}{19}
% started from zero:
% (let ((year 2013) (month 7)) (+ (* (- year 2005) 12) month -1))
\newcommand{\debmtgnumber}{115}

% tikz picture の為のマクロ設定
\usepackage[dvipdfmx]{graphicx}
\usepackage{tikz}

\begin{document}

\begin{titlepage}
\thispagestyle{empty}
% タイトルページ:編集必要な部分は最初のマクロに飛ばすこと

\vspace*{-2cm}
第\debmtgnumber{}回 東京エリア Debian 勉強会資料\\
\hspace*{-2cm}
\includegraphics{image2012-natsu/dotdeb.pdf}\\
\hfill{}\debmtgyear{}年\debmtgmonth{}月\debmtgdate{}日

% ここはアップデートすること
% 全角文字にしないとフォントのサイズが合わないので注意
\rotatebox{10}{\fontsize{30}{30} {\gt 特集：Jenkins での Debian パッケージ自動化}}\\

\vspace*{-2cm}
\hfill{}\includegraphics[height=6cm]{image200502/openlogo-nd.eps}
\end{titlepage}

\newpage

\begin{minipage}[b]{0.2\hsize}
 \definecolor{titleback}{gray}{0.9}
 \colorbox{titleback}{\rotatebox{90}{\fontsize{80}{80} {\gt デビアン勉強会} }}
\end{minipage}
\begin{minipage}[b]{0.8\hsize}
\hrule
\vspace{2mm}
\hrule
\begin{multicols}{2}
\tableofcontents
\end{multicols}
\vspace{2mm}
\hrule
\end{minipage}

\dancersection{事前課題}{野島 貴英}

今回の事前課題は以下です:
\begin{enumerate}
 \item 本日、何の作業をやるかを宣言ください。
\end{enumerate}
この課題に対して提出いただいた内容は以下です。
\begin{multicols}{2}
{\small
 \begin{prework}{ koedoyoshida }
未定
\end{prework}

\begin{prework}{ dictoss(杉本　典充) }
パッケージ作成に関する作業は、\url{http://wnpp.debian.net/}に報告が出ているパッケージ化、またはパッチつくりをやろうと思います。具体的にはmunin、pygtkplot、pylzma、wx3.0-doc、hex2binが候補です。
\end{prework}

\begin{prework}{ 工藤寿晴 }
自分がVineで管理しているパッケージから、何かDebian用のパッケージを作ります。
（フォントかCUIのソフトウェアを検討中）
\end{prework}

\begin{prework}{ 吉野(yy\_y\_ja\_jp) }

manpages-ja
\end{prework}

\begin{prework}{ zinrai }

パッケージ作成

pulledpork\\
\url{https://code.google.com/p/pulledpork/}
\end{prework}

\begin{prework}{ wbcchsyn }

GitHub で公開中の自作プログラムをパッケージする予定
(先月の勉強会で発表した物です。)\\
\url{https://github.com/wbcchsyn/dVault}

時間があまったら、プログラムの機能実装を行います。
\end{prework}

\begin{prework}{ 鈴木崇文 }

Cocotのパッケージングをやります。
\end{prework}

\begin{prework}{ 野島　貴英 }

パッケージング道場参加します。
\begin{itemize}
\item xmrisのパッケージング
\item kanatestのバグつぶし
\item Softetherのパッケージ化ヒントが得られるまで 
\end{itemize}
となってます。
\end{prework}

\begin{prework}{ 岩松 信洋 }

パッケージアップロードします。
\end{prework}

\begin{prework}{ まえだこうへい }
\begin{itemize}
\item ToDo潰し\\
\url{http://d.palmtb.net/pages/about.html#todo-from-maintainer-dashboard}
\item pendingのままのSerf関連のパッケージ化
\end{itemize}
を行う予定。

\end{prework}

}
\end{multicols}

\dancersection{Debian Trivia Quiz}{野島 貴英}

ところで、みなさん Debian 関連の話題においついていますか？Debian関連の話
題はメーリングリストをよんでいると追跡できます。ただよんでいるだけではは
りあいがないので、理解度のテストをします。特に一人だけでは意味がわからな
いところもあるかも知れません。みんなで一緒に読んでみましょう。

今回の出題範囲は\url{debian-devel-announce@lists.debian.org} や \url{debian-devel@lists.debian.org}に投稿された
内容などからです。

\begin{multicols}{2}
%; whizzy-master ../debianmeetingresume201311.tex
% 以上の設定をしているため、このファイルで M-x whizzytex すると、whizzytexが利用できます。
%

\santaku
{libcが変更になりました。変更前と変更後の組み合わせで正しいのはどれ？}
{eglibcからglibcへ変更}
{glibcからeglibcへ変更}
{uClibcからglibcへ変更}
{A}
{eglibcが採択されてから5年が経過し、glibcに戻ってきた状況です。}

\santaku
{ネットワークはTorを使い、ファイル/メール/メッセンジャーソフトは暗号化対応という特徴を持つdebianベースのLive OSの名前は何でしょう？}
{Plan9}
{Tails}
{GNU/Hurd}
{B}
{TrailsはThe Amnesic Incognito Live Systemの略です。利用者の出来る限りの匿名性を保つ事を目的として作られたDebianベースのlive osとなります。}

\santaku
{Debian 6 LTS開発の為の寄付を受け付ける会社名は？}
{Canonical}
{SquareEnix}
{Freexian}
{C}
{Freexian社は、Debianの公式開発者のRapha\"el Hertzogにより設立された、
Debianの有償サポートをするための会社です。}

\santaku
{jessieのデフォルトのgccのバージョンは何になった？}
{4.9}
{4.6}
{2.6}
{A}
{各パッケージはgcc-4.9でもコンパイルが通るように修正が必要となります。}

\santaku
{jessieのfreezeは11月5日です（注：リリースじゃないよ。）ここで、スケジュールで新規のtransitionが出来なくなる日はいつでしょう？}
{10月5日}
{9月5日}
{8月5日}
{B}
{9月5日以降は新規のtransitionが出来なくなります。新しいパッケージをjessieに合わせてリリースしようと目論む方々は期日にご注意下さい。}





\end{multicols}

\dancersection{最近のDebian関連のミーティング報告}{野島 貴英}

\subsection{東京エリアDebian勉強会114回目報告}

 東京エリアDebian勉強会114回目は(株)スクウェア・エニックスさんで開催されました。
 
 9人の参加者がありました。また、wbcchsynさんにて、GPGの秘密鍵を堅牢にかつセキュアに
管理する事を検討した成果についての発表がありました。今回は、その成果の１つとして、
堅牢化についてフォーカスし、RAIDの考え方を用いて複数の記憶メディアにデータを
分散配置して復元するソフトウェアをpythonで実装した事についての発表でした。また、
こちらについて様々に意見交換が参加者の間で行われました。

　発表者自身のアイデアに基づくDebianに関係した成果物の発表は勉強会としても
大変価値があると考えております。引き続き、発表者の方々によるユニークなアイデアと
成果の発表が期待されます。

% % (query-replace-regexp "<.*?>" "")
% % (query-replace-regexp "^[	 ]\+" "")

%-------------------------------------------------------------------------------
\dancersection{JenkinsでのDebianパッケージ自動化}{前田耕平}
%-------------------------------------------------------------------------------
\index{building-packages-automatically-on-jenkins}

昨年の大統一Debian勉強会2013で、「とある企業でのDebianシステムの使い方。」というセッション\footnote{\url{http://gum.debian.or.jp/2013/session/437.html}}の中で、社内向けのDebianパッケージの作成と、repreproを使ったパッケージの配布方法についての話をしました。今回はその続編です。

\subsection{背景}

管理するパッケージが少ない場合は、手作業でのビルドでも間に合っていましたが、管理するパッケージ数が増え、対象ディストリビューションが増え(Precise, Wheezy, Trusty)、一方メンテナンスする人を増やすのは時間が掛かる、という状況でした。

そこで、別件\footnote{JavaやPythonのプロジェクトのビルドやテストの実行目的。}でJenkinsを用意したこともあったので、ソースパッケージを作成するプロセス以外は、全てJenkinsで実行して自動化することにしました。

\subsection{自動化するプロセス}

パッケージのビルドからローカルアーカイブへの登録までに行うプロセスは図1のようになります。
1aは公式パッケージのソースパッケージをdgetなどで取得する場合、1bは独自にソースパッケージを作成する場合です。ビルドする対象のソースパッケージによって実際の手段自体は異なってきますが、実行するプロセス自体は図の通り変わりません。

\begin{figure}[H]
\begin{center}
\begin{tikzpicture} [x=1em, y=1ex]
  \draw (-33, -1) rectangle (-17, -5);
  \draw (-25, -3) node {1a. ソースパッケージ取得};
  \draw[->] (-25, -5) -- (-17, -8);

  \draw (-5, -1) rectangle (11, -5);
  \draw (3, -3) node {1b. ソースパッケージ作成};
  \draw[->] (3, -5) -- (-5, -8);

  \draw (-23, -8) rectangle (1, -12);
  \draw (-11, -10) node {2. パッケージビルド};
  \draw[->] (-11, -12) -- (-11, -15);

  \draw (-23, -15) rectangle (1, -19);
  \draw (-11, -17) node {3. パッケージチェック};
  \draw[->] (-11, -19) -- (-11, -22);

  \draw (-23, -22) rectangle (1, -26);
  \draw (-11, -24) node {4. インストールテスト};
  \draw[->] (-11, -26) -- (-11, -29);

  \draw (-23, -29) rectangle (1, -33);
  \draw (-11, -31) node {5. ソースパッケージへの署名};
  \draw[->] (-11, -33) -- (-11, -36);

  \draw (-23, -36) rectangle (1, -40);
  \draw (-11, -38) node {6. ローカルアーカイブへの登録};
\end{tikzpicture}
\caption{パッケージビルド〜ローカルアーカイブへの登録までのプロセス}
\end{center}
\end{figure}

各プロセスの実行内容は下記のようになります。

\begin{figure}[H]
\begin{center}
\begin{tikzpicture} [x=1em, y=1ex]
  \draw (-33, -1) rectangle (-17, -5);
  \draw (-25, -3) node {1a. dgetで公式アーカイブから取得};
  \draw[->] (-25, -5) -- (-25, -8);

  \draw (-5, -1) rectangle (11, -5);
  \draw (3, -3) node {1b. git-buildpackageでGitで管理};
  \draw[->] (3, -5) -- (3, -8);

  \draw (-33, -8) rectangle (-17, -12);
  \draw (-25, -10) node {2a. cowbuilder/pbuilderでビルド};
  \draw[->] (-25, -12) -- (-17, -15);

  \draw (-5, -8) rectangle (11, -12);
  \draw (3, -10) node {2b. git buildpackageでビルド};
  \draw[->] (3, -12) -- (-5, -15);

  \draw (-23, -15) rectangle (1, -19);
  \draw (-11, -17) node {3. ビルド時のhook scriptとしてlintianでチェック};
  \draw[->] (-11, -19) -- (-11, -22);

  \draw (-23, -22) rectangle (1, -26);
  \draw (-11, -24) node {4. piupartsでインストール＆アンインストールテスト};
  \draw[->] (-11, -26) -- (-11, -29);

  \draw (-23, -29) rectangle (1, -33);
  \draw (-11, -31) node {5. debsignで署名};
  \draw[->] (-11, -33) -- (-11, -36);

  \draw (-29, -36) rectangle (7, -43);
  \draw (-11, -38) node {6-1. dputでincomingディレクトリにpush};
  \draw (-11, -41) node {6-2. dputのpost\_upload\_commandとしてreprepro processincommingを実行};
\end{tikzpicture}
\caption{各プロセスの実行内容}
\end{center}
\end{figure}

これらをJenkinsでジョブとして実行するためには、
対話形式の処理をなくし、リモートホスト上の処理も実行可能な状態で、
一連の処理をJenkinsのジョブスクリプトを作成する必要があります。

\subsubsection{Jenkinsで実行する上での課題}

Jenkinsで実行する上で問題となるのは、GPGとSSH公開鍵のパスフレーズの入力と、
repreproでのアーカイブの登録です。

repreproでのアーカイブ登録は\texttt{reprepro processincoming}コマンドで行います。Jenkinsと別ホストでrepreproを実行する場合、dput.cfの\texttt{post\_upload\_command}を使えば解決できます。
例えば、Ubuntu Trusty用のローカルアーカイブにdputする場合は、下記のような設定をdput.cfにしておきます。

\begin{commandline}
[trusty]
fqdn                    = repo.example.org
method                  = scp
incoming                = /var/lib/debpkg-custom/ubuntu/incoming/
login                   = repoadm
post_upload_command     = ssh repoadm@repo.example.org /usr/bin/reprepro -b /var/lib/debpkg-custom/ubuntu \
--keepunreferencedfiles processincoming trusty #実際は一行
\end{commandline}

公開鍵の入力は、次の3箇所で必要になります。\footnote{Jenkinsとrepreproを同じホストで実行する場合(つまり一台で完結させる場合)にはSSH公開鍵のパスフレーズ入力は特に考慮することはありません。}

\begin{itemize}
  \item GPG
    \begin{itemize}
    \item[5.] debsignでの署名
    \item[6-2.] dputのpost\_upload\_commandとしてreprepro processincommingを実行
    \end{itemize}
  \item SSH
    \begin{itemize}
    \item[6-1.] dputでincomingディレクトリにpush
    \end{itemize}
\end{itemize}

一方、GPGについては既存環境で使っている鍵をそのまま利用する関係上、避けて通れません。
Jenkinsのジョブスクリプトで

\begin{commandline}
$ echo -e 'passphrase\npassphrase\n' | debsign some.changes
$ echo -e 'passphrase\npassphrase\n' | dput some.changes
\end{commandline}

と実行しても標準入力からはパスフレーズは渡せません。しかし、gpgコマンド自体には\texttt{--no-tty --batch}オプションを使って、対話形式でなくてもパスフレーズを入力することができます。そのため、方法としては、
\begin{itemize}
  \item debsign, reprepro processincommingをTTYなしで対話形式でなくてもGPGパスフレーズを入力できるようにする
  \item expectを使ってプロンプトで判断してパスフレーズを入力させる
\end{itemize}

の2つの方法が考えられます。
で、先に着手したdebsignについては、expectで処理するのは負けだと思ったので、別途pydebsignというPythonライブラリを実装しました\footnote{Git: \url{https://github.com/mkouhei/pydebsign}, PyPI: \url{https://pypi.python.org/pypi/pydebsign}}。しかし、repreproについては時間の都合上、結局Python版のexpectのpexpectを使いました。下記のようになります。\footnote{debsignのこれで良かったですね。。。}

\begin{commandline}
def upload_files(codename, changes_path, reprepro_passphrase):
    """ uploading to local archive incoming directory of reprepro """
    os.environ['LANG'] = 'C'
    command = "dput %s %s" % (codename, changes_path)
    exp = pexpect.spawn(command, maxread=4000)
    exp.logfile_read = sys.stdout
    exp.expect('Please enter passphrase:')
    exp.sendline(reprepro_passphrase)
    exp.expect('Please enter passphrase:')
    exp.sendline(reprepro_passphrase)
    exp.expect(pexpect.EOF)
    exp.close()
    return exp.status
\end{commandline}

\subsubsection{Jenkinsで必要なプラグイン}

パスフレーズは渡せるようになりましたが、Jenkinsのジョブを実行するとパスフレーズが丸見えです。そのため、"Mask Passwords Plugin"\footnote{\url{https://wiki.jenkins-ci.org/display/JENKINS/Mask+Passwords+Plugin}}というプラグインを使います。
他に必要なプラグインは下記の通りです。また、Jenkinsのジョブスクリプト管理用と、git-buildpackage用と複数のGitリポジトリを扱えるようにするため、"Git Plugin"\footnote{\url{https://wiki.jenkins-ci.org/display/JENKINS/Git+Plugin}}と"Multiple SCMs Plugin"\footnote{\url{https://wiki.jenkins-ci.org/display/JENKINS/Multiple+SCMs+Plugin}}も必要です。これらのプラグインは後述のJenkins環境構築後に、Jenkinsの管理画面からインストールします。

\subsubsection{その他の考慮点}
\begin{itemize}
\item Jenkinsで自動ビルド＆ローカルアーカイブへ登録したパッケージに依存するパッケージを作成することもあります。そのため、後述のpbuilderとcowbuilderのイメージのAPT lineには、ローカルアーカイブのURLを登録し、ローカルアーカイブの公開鍵をapt-keyで登録しておく必要があります。
\item バックポートする場合、orig.tar.{gz,xz}が.changesファイルのリストにはデフォルトでは登録されず、repreproで登録する際にエラーになります。そのため、debbuildopts='-sa'を指定する必要もあります。
\item pbuilderやcowbuilderでパッケージをビルドすると、デフォルトでは/var/cache/pbuilder/resultsディレクトリ以下に作成されます。このままでは、Jenkinsのジョブごとに処理するには不都合です。pbuilderrcでBUILDRESULTに、Jenkinsのworkspace以下のパスを指定するようにします。
\end{itemize}

以上の点を考慮し、PythonでJenkinsジョブ用のスクリプトを作成しました。\footnote{デモのみ。}
都合により、Golang版も作成しました。\footnote{\url{https://github.com/mkouhei/godebbuild}}
これらは、前述のhook scriptで行うlintianでのテストを、dputの\texttt{-l}オプションで実行できるようにしています。\footnote{backportなどの場合、lintianのWarning, Errorを修正するためのパッケージの変更せず、許容してしまいたいこともあるため。}

\subsection{Jenkins及びreprepro環境の構築}

Jenkinsで自動実行するための環境として、下記の準備が必要になります。

\begin{itemize}
  \item Jenkinsのインストール＆設定
  \item Jenkins用ユーザへのsudo設定
    \begin{itemize}
      \item /usr/sbin/pbuilder
      \item /usr/sbin/cowbuilder
      \item /usr/bin/git-buildpackage
      \item /usr/bin/piuparts
      \item /bin/chown
    \end{itemize}
  \item pbuilder and cowbuilderイメージの作成
    \begin{itemize}
      \item ビルド対象のディストリビューションの数分だけ作成
    \end{itemize}
  \item debsignおよびreprepro用のGPG鍵の作成
  \item repreproへのupload用のSSH公開鍵の作成\footnote{Jenkinsとrepreproが別ホストの場合}
  \item dput.cfの設定
  \item rererproのインストール＆設定
  \item reprepro管理poolの公開設定
  \item ローカルアーカイブ用のGPG公開鍵のexport
\end{itemize}

このうち、pbuilderとcowbuilderイメージの作成の一部と、GPG鍵の作成、SSH公開鍵の作成については、手動で行う必要がありますが、他はAnsibleで自動化しました。\footnote{デモ参照。Playbookは、\url{https://github.com/mkouhei/playbook-tokyodebian115}に用意してあります。}

\subsection{まとめ}

ソースパッケージの作成とJenkinsのジョブの作成以外は自動化できました。
これにより、Debianパッケージの作成をやったことが無い人でもバックポートなどはできるようになりました。
しかし、実際にはジョブを実行するとパッケージの依存関係や、lintian、piupartsなどのエラーが発生します。そのため、Debianパッケージ作成について知らなければ、それらのエラーに対応できないことが予想されます。

また、ソースパッケージの作成は相変わらず、属人化しています。\footnote{つまり、職場では自分しか作成できない。}
なので、今回のDebianパッケージング道場でパッケージ作成ができる人が増えるとイイですね。\footnote{原稿作成時点では職場の同僚の参加者は未定。}

\printindex

\cleartooddpage

\vspace*{15cm}
\hrule
\vspace{2mm}
\includegraphics[width=2cm]{image200502/openlogo-nd.eps}
\noindent \Large \bf Debian 勉強会資料\\
\noindent \normalfont \debmtgyear{}年\debmtgmonth{}月\debmtgdate{}日 \hspace{5mm}  初版第1刷発行\\
\noindent \normalfont 東京エリア Debian 勉強会 （編集・印刷・発行）\\
\hrule

\end{document}
